% ---- pg-263 （书页 251） ----
\[
T^2=[-0.167,-0.066,0.200]\left(\frac{1}{30}+\frac{1}{30}\right)^{-1}\begin{bmatrix}0.719&-0.268&-0.125\\-0.268&1.101&-0.751\\-0.125&-0.751&1.058\end{bmatrix}^{-1}\begin{bmatrix}-0.167\\-0.066\\0.200\end{bmatrix}
\]
\[
=15\times0.067=1.005
\]
此外, 对 $\alpha=0.05$, $c^2=[(30+30-2)(4-1)/(30+30-4)]F_{3,56}(0.05)=3.11\times2.8=8.7$. 因为 $T^2=1.005<8.7$, 我们得出结论: 两个总体的平行轮廓假设是可以接受的. 由\hyperref[fig:6.5]{图 6.5} 可以看出, 这一结论并不奇怪.

在接受平行轮廓的假设之后, 我们要检验轮廓的相合性, 即检验 $H_{02}: \mathbf{1}'\boldsymbol{\mu}_1=\mathbf{1}'\boldsymbol{\mu}_2$(轮廓重合). 为此需要完成下列计算:

\noindent $\bar{\mathbf{x}}_1-\bar{\mathbf{x}}_2$ 中元素之和 $=\mathbf{1}'(\bar{\mathbf{x}}_1-\bar{\mathbf{x}}_2)=0.367$

\noindent $\mathbf{S}_p$ 中元素之和 $=\mathbf{1}'\mathbf{S}_p\mathbf{1}=4.207$

利用\hyperref[eq:6-74]{式(6-74)}得到
\[
T^2=\left[\frac{0.367}{\sqrt{\left(\frac{1}{30}+\frac{1}{30}\right)\times4.027}}\right]^2=0.501
\]
对 $\alpha=0.05$ 有 $F_{1,58}(0.05)=4.0$ 和 $T^2=0.501<F_{1,58}(0.05)=4.0$, 所以我们不能拒绝轮廓的相合性假设, 即对所考虑的四个问题, 男性与女性的回答并没有差异.

现在可以检验这两个轮廓的水平性. 但由于问题 1 和 2 的衡量尺度是 1$\sim$8 分, 而问题 3 与 4 的衡量尺度则是 1$\sim$5 分, 所以在本例中这种检验是没有意义的. 尺度的不相容性使水平轮廓的检验失去意义, 这说明了给出一个完整的轮廓分析, 需要有相似的测量值.\hfill$\blacksquare$

当样本容量较小时, 轮廓分析依赖于正态性假定. 用第 4 章讨论的方法, 可用原始观测值 $\mathbf{x}_{ij}$ 或对照观测值 $\mathbf{C}\mathbf{x}_{ij}$ 对这个假定进行检验.

对于两个以上的总体, 其轮廓分析可仿照两总体的做法去完成. 事实上, 就连那些比较的度量也可从上面所讨论的度量通过类比得出(见文献[13]和[18]).

\section{重复测量设计和生长曲线}  % 原书 6.9
\phantomsection\label{sec:6.9}

如前所述, 重复测量是针对在不同的时间和地点, 对同一对象(个体)观测其同一特征这种情况来说的.

(a) 一个个体的观测值可与各种不同的处理相对应. 比如在例 6.2 中, 对每一只狗, 可在 $2\times2$ 种处理组合下分别测量其心跳的间隔时间. 当同一个体的诸响应相关时, 还需比较它们的处理.

(b) 一种处理可用于每一个个体, 并在一个时间段内观测一个特征. 例如, 我们可以测量一只刚出生的幼狗的体重, 而后每过一月测量一次, 这样就能描绘一只典型的狗生长轨迹的曲线. 我们在本书中称这条曲线为生长曲线.

当某些个体接受一种处理而另一些个体接受另一种处理时, 就需要对这两种处理的生长曲线进行比较.

为说明普特霍夫与罗伊(文献[21])所提出的生长曲线模型, 我们考虑老年妇女主尺骨中的钙含量问题. 对于控制组, 除了测试开始时记录的数据之外, \hyperref[tab:6.5]{表 6.5} 中还给出了经过 1 年、2 年和 3 年后钙含量数据. 对同一个体(即测试对象), 由光子吸收测量术获得的测试数据是相关的, 但不同个体间的测试数据应相互独立. 这种模型假定每一个个体均有相同的协方差矩
